\section{Related Work}\label{sec:related}
This paper extends the work of \citet{oopsla11} with variance for
polymorphic types and a dynamic dispatch algorithm.
Allen et al. proposed rules for defining overloaded functions that ensure type safety
under symmetric multiple dispatch in an object-oriented language with
multiple inheritance, and showed how to check the rules without requiring
the entire type hierarchy.
They proposed the overloading rules in the context of a set of class declarations
and function declarations without defining a language with complete
static and dynamic semantics.  Therefore, their work does not guarantee
that the static overloading rules and the way a language chooses a method
to invoke at run time match correctly.
% They did not provide a complete static and dynamic semantics; therefore, their work does not guarantee
% that the static overloading rules and the way a language chooses a method
% to invoke at run time match correctly.
While FGFV does not support features like modularity, nominal exclusion,
and closed types in their work, they can be added in an obvious way.
As the related work section of~\citet{oopsla11} discusses in detail,
most earlier systems like \textit{Dubious}~\cite{dubious99,dubious02,dubious03},
MultiJava~\cite{multijava06}, and F(E{\footnotesize{ML}})~\cite{FEML06} supported
overloading with symmetric multiple dynamic dispatch with multiple inheritance,
but without support for polymorphic methods or types.  While $\text{ML}_\le$~\cite{bm97}
integrates polymorphism and symmetric multiple dynamic dispatch, it lacks
multiple inheritance.

ParaSail~\cite{parasail11,parasailweb} is a parallel programming language designed to support
inherently safe and secure, highly parallel applications.
It also constrains the scope of type parameters,
prohibiting a bound on a type parameter from referring to type parameter
bindings to its left, so that dynamic type inference can be fast.
However, ParaSail does not support covariant and contravariant
type parameters, and it does not provide any formal specification.

Julia~\cite{julia15,julia17} is a dynamic programming language with symmetric
multiple dispatch. It provides a type declaration and type annotation
sublanguage, but does not have a static type system. Thus, it does not
check the validity of overloaded methods, and only uses the
subtype relation to find the most specific method among
overloaded methods at run time.

Various languages provide different subsets of the FGFV features.
Scala~\cite{scala} supports polymorphic types
with variance, but it does not support symmetric multiple dispatch.
FHJ~\cite{fhj} proposed a dispatch algorithm, which uses static types at run time 
to resolve overloading conflicts from unrelated methods that happen to have
the same name and signature due to multiple inheritance.
\citet{castagna15} presented a statically-typed calculus with dynamic
(explicit) dispatch through pattern-matching, polymorphism, recursive
types, semantic subtyping, and union, intersection, and negation type
connectives.  Their local type inference algorithm computes sets of type
substitutions to insert in appropriate places in expressions.  While
their type-case expression and intersection type behave similarly to
method overloading, they do not have a similar notion of overloading
rules.  Since the type-case expression executes the first matching
expression, it does not introduce any ambiguity.
For OCaml, Garrigue~\citeyearpar{garrigue01,garrigue02} introduced a type
inference system for structural polymorphism such as polymorphic records,
polymorphic variants, and accurate matching.  Type constraints are
described by constraint domains that are defined separately from its
type system, and the constraints contain some notion of lower and upper
bounds, which respectively represent required and available labels or cases.
He also extended his work with CoQ mechanization \cite{garrigue10,garrigue15}.
The work of Castagna and Garrigue may serve as a possible starting point
for finding the necessary substitution $\sigma$ to satisfy the rule
\rncite{E-Sub} for existential subtyping during static analysis.
We believe that our work opens the gate to adding either variance
annotation in the work of Castagna or a form of type-case in the work of
Garrigue.
